\hypertarget{implementation-how-the-system-works}{%
\chapter{Implementation: how the system
works}\label{implementation-how-the-system-works}}

\begin{draftnote}

\textbf{Draft status, 03 November 2026.} Written from the implemented
code and from \texttt{docs/PROJECT\_STATUS.md}. Every constant quoted
below is named in the source; \textbf{check each against the file cited
beside it before submission}, because a constant quoted wrongly in a
report is the easiest thing for an examiner to catch and the hardest to
recover from. Figures 1 to 5 referred to here are the committed diagrams
in \texttt{docs/diagrams/}.

\end{draftnote}

\begin{center}\rule{0.5\linewidth}{0.5pt}\end{center}

\hypertarget{the-stack-and-why-each-part-is-where-it-is}{%
\section{The stack, and why each part is where it
is}\label{the-stack-and-why-each-part-is-where-it-is}}

\begin{longtable}[]{@{}
  >{\raggedright\arraybackslash}p{(\columnwidth - 4\tabcolsep) * \real{0.3333}}
  >{\raggedright\arraybackslash}p{(\columnwidth - 4\tabcolsep) * \real{0.3333}}
  >{\raggedright\arraybackslash}p{(\columnwidth - 4\tabcolsep) * \real{0.3333}}@{}}
\toprule\noalign{}
\begin{minipage}[b]{\linewidth}\raggedright
tier
\end{minipage} & \begin{minipage}[b]{\linewidth}\raggedright
technology
\end{minipage} & \begin{minipage}[b]{\linewidth}\raggedright
reason
\end{minipage} \\
\midrule\noalign{}
\endhead
\bottomrule\noalign{}
\endlastfoot
Web client & React 19, Vite & the examiner-facing demonstration; holds
no credentials \\
API & Node, Express 5, TypeScript, Mongoose & the only tier that holds
credentials or reaches a store \\
Detector & Python, FastAPI, tree-sitter, APTED & the parsing and tree
libraries that exist are Python ones \\
Documents & MongoDB & the detection result is a nested document, not a
set of rows \\
Queue & Redis & detection is slow and must not block an upload \\
Files & MinIO-compatible object store & submitted source is a blob, not
a column \\
\end{longtable}

The split between API and detector is the one design decision worth
defending in a viva, because it looks like unnecessary complexity for a
student project. It is not. The parsing ecosystem this project needs ---
tree-sitter grammars and an APTED implementation --- is Python. The web
application and its data model are far more naturally expressed in
TypeScript against Mongoose. Rather than compromise one to accommodate
the other, the system makes the detector a stateless HTTP service that
parses and compares, and keeps every piece of state, every credential
and every authorisation decision in the API.

The worker is deliberately \emph{not} a separate codebase. It is the API
started on a different entry point, so a queued detection job and a
synchronous request share one set of services. Two implementations of
the same scoring logic would drift, and the drift would appear as a
result that cannot be reproduced.

\hypertarget{layer-1-structural-similarity-in-working-order}{%
\section{Layer 1 --- structural similarity, in working
order}\label{layer-1-structural-similarity-in-working-order}}

Layer 1 answers \emph{does this submission resemble another submission?}
Its chain runs through four stages, and the stages matter because two of
them operate on different representations of the same file.

\hypertarget{parsing}{%
\subsection{Parsing}\label{parsing}}

\texttt{app/parsing.py} exposes
\texttt{parse\_source(language,\ source)}, returning a
\texttt{ParseResult} carrying \texttt{ok}, \texttt{node\_count},
\texttt{error} and \texttt{root}. One detail is load-bearing and caused
a real defect during the measurement work: \textbf{the root is returned
even when the parse fails.} tree-sitter \cite{treesitter} is an
error-recovering parser, so a syntactically broken file still yields a
tree. Any caller that tests \texttt{root\ is\ None} rather than
\texttt{ok} will therefore proceed on a tree built from broken source.
Callers test \texttt{ok}.

\hypertarget{normalisation}{%
\subsection{Normalisation}\label{normalisation}}

\texttt{app/normalise.py} turns the raw tree-sitter tree into a
\texttt{TNode} tree via \texttt{normalise(node)}. Normalisation is where
the resistance to obfuscation comes from: identifier names, literal
values and formatting are discarded, and each node carries a
\texttt{label} describing its syntactic kind rather than its text. This
is what makes renaming a variable invisible to Layer 1.

\textbf{The two representations are not interchangeable, and the report
should say so plainly}, because it is the clearest illustration of what
each layer is looking at:

\begin{itemize}
\tightlist
\item
  Layer 1 consumes the \textbf{normalised} tree, because structure is
  what survives obfuscation.
\item
  Layer 2 consumes the \textbf{raw} tree, because style lives precisely
  in the things normalisation throws away --- blank lines, comments,
  line lengths.
\end{itemize}

The same source file is therefore read twice, in two different ways, for
two different questions. \texttt{extract\_units} and \texttt{size} take
the normalised tree; \texttt{extract\_features} takes the raw root and
the source bytes.

\hypertarget{unit-extraction}{%
\subsection{Unit extraction}\label{unit-extraction}}

\texttt{app/units.py} provides \texttt{extract\_units(tree,\ language)},
which splits a normalised tree into comparable top-level units ---
functions, methods, classes --- rather than treating the file as one
monolithic tree. Comparing whole files would make a long honest file
resemble another long honest file simply by virtue of both being long.
Comparing unit against unit allows the system to report \emph{which}
parts matched, which is what a reviewer needs.

\hypertarget{prefilter-the-step-that-makes-it-affordable}{%
\subsection{Prefilter: the step that makes it
affordable}\label{prefilter-the-step-that-makes-it-affordable}}

\texttt{app/prefilter.py} exists because tree edit distance is
expensive. It builds a \texttt{Counter} of node labels per tree ---
\texttt{label\_counts} --- and scores candidate pairs by bag overlap
with \texttt{quick\_similarity}, which is cheap and approximate and
never final. \texttt{rank} then nominates the most promising candidates.

Two constants from \texttt{app/main.py} govern the cost:

\begin{itemize}
\tightlist
\item
  \texttt{MAX\_CANDIDATES\ =\ 50} --- the most candidates the endpoint
  will accept at all.
\item
  \texttt{MAX\_DEEP\_CANDIDATES\ =\ 10} --- of those, how many earn the
  expensive treatment.
\end{itemize}

So a submission is compared in full against at most ten candidates
however large the cohort is. This bounds one analysis irrespective of
class size, and it is also a \textbf{fidelity requirement on any
measurement of Layer 1}: a study that compares exhaustively is measuring
a system that was never deployed, and will credit the layer with matches
the API cannot find. That error was made and corrected during this
project's own ablation work, and Chapter 7 records it.

\texttt{shortlist\_pairs} nominates pairs from \emph{both} directions.
Shortlisting in one direction starves units nobody happens to pick: if
every unit of A nominates the same five units of B, the remainder of B
is never offered to anyone and their partners in A go unmatched.

\hypertarget{comparison-and-scoring}{%
\subsection{Comparison and scoring}\label{comparison-and-scoring}}

\texttt{compare\_unit\_sets(units\_a,\ units\_b,\ size\_a,\ size\_b)}
returns a \texttt{FileComparison} carrying the similarity, the matched
unit pairs, coverage on each side, and counts of pairs compared, pruned
and skipped. The underlying distance is APTED \cite{pawlik2016}, with
RTED \cite{pawlik2015} as the algorithmic predecessor the authors ask to
be cited alongside.

Reported evidence is filtered: \texttt{MIN\_SPAN\_NODES\ =\ 10}
suppresses matches too small to mean anything, because two one-line
functions resemble each other by virtue of being one line long. The
filter applies to the \emph{spans shown to a reviewer}, not to the
similarity figure itself --- a distinction worth knowing before
answering a question about it.

\hypertarget{layer-2-behavioural-fingerprinting-in-working-order}{%
\section{Layer 2 --- behavioural fingerprinting, in working
order}\label{layer-2-behavioural-fingerprinting-in-working-order}}

Layer 2 answers a different question: \emph{does this submission
resemble its own author's previous work?} It is the only layer that can
see outsourced work, because it needs no second copy anywhere.

\hypertarget{features}{%
\subsection{Features}\label{features}}

The detector's \texttt{/features} endpoint extracts \textbf{ten}
features via \texttt{app/features.py}: \texttt{avg\_line\_length},
\texttt{blank\_line\_ratio}, \texttt{comment\_density},
\texttt{avg\_identifier\_length},
\texttt{underscore\_identifier\_ratio}, \texttt{avg\_function\_lines},
\texttt{functions\_per\_100\_lines}, \texttt{max\_block\_depth},
\texttt{avg\_params\_per\_function} and \texttt{for\_loop\_ratio}. All
ten are stored on the baseline.

\textbf{Four are scored.} \texttt{SCORED\_FEATURES} in
\texttt{services/behavioural.ts} is \texttt{blank\_line\_ratio},
\texttt{avg\_line\_length}, \texttt{max\_block\_depth},
\texttt{comment\_density}. Storing all ten keeps the choice reversible.

The justification for choosing those four is \textbf{retracted}, and the
report must say so rather than repeat it. The four were selected from a
60-file sample; a later measurement over 204 Python files across 15
repositories found that the earlier ranking does not replicate, that
\texttt{blank\_line\_ratio} --- which had driven the choice --- ranks
ninth of ten in Python, and that using every applicable feature beats
the fixed four in both languages. Chapter 7 carries the figures.

\hypertarget{anchors-and-the-baseline}{%
\subsection{Anchors and the baseline}\label{anchors-and-the-baseline}}

\texttt{services/anchors.ts} decides which past submissions may describe
a student's style, and \texttt{services/baseline.ts} builds the
per-student, per-language profile from them. Eligibility is the first of
the five integrity mitigations below.

\hypertarget{the-score-and-what-its-two-halves-mean}{%
\subsection{The score, and what its two halves
mean}\label{the-score-and-what-its-two-halves-mean}}

\begin{verbatim}
z        = (value - mean) / stdDev
stdDev   = max(cohortSpread, VARIANCE_FLOOR[feature])
score    = min(1, rootMeanSquare(z) / Z_SATURATION)      Z_SATURATION = 3
\end{verbatim}

The asymmetry here is the project's central idea and the thing most
likely to be misread, including in its own earlier documentation:

\begin{itemize}
\tightlist
\item
  the \textbf{numerator} is the student's own trust-weighted mean across
  their anchors;
\item
  the \textbf{denominator} is the \textbf{cohort's} spread, floored.
\end{itemize}

A z-score therefore means \emph{how far this submission sits from this
student's usual, measured in units of how much students in this class
differ from one another}. Both halves are load-bearing. An earlier
version of this project's documentation claimed the comparison was
``never against the cohort'', which is false and does not survive an
examiner opening the file.

The cohort denominator is a measured choice rather than a convenience: a
student's own spread, taken from three or four anchors, is too unstable
to divide by. A student's mean is personal; the spread is closer to a
property of the feature. The variance floor exists because a cohort can
produce zero spread on a coarse feature, and dividing by zero gives
infinity.

\textbf{The AUC figures previously quoted here are retracted} --- 0.766
for a student's own spread against 0.909 for the cohort's. No
measurement script computes an AUC; see Chapter 7 \S{}7.4. The cohort
denominator rests on the stability argument above, which is reasoning
rather than measurement.

\hypertarget{the-five-baseline-integrity-mitigations}{%
\subsection{The five baseline-integrity
mitigations}\label{the-five-baseline-integrity-mitigations}}

Layer 2's failure mode is not a false negative; it is a wrong baseline
accusing an honest student. Five mitigations exist, and they are a
genuine contribution of the project rather than boilerplate:

\begin{enumerate}
\def\labelenumi{\arabic{enumi}.}
\tightlist
\item
  \textbf{Anchor provenance.} Only invigilated work, or work a student
  nominates and faculty accept, may describe their style.
\item
  \textbf{Minimum anchor count.} Below \texttt{minAnchorsForBaseline}
  (3) the layer declines to score rather than guessing.
\item
  \textbf{Baseline confidence.} A thin or inconsistent baseline
  attenuates the layer's weight (\texttt{effectiveW2}) instead of being
  trusted at full strength.
\item
  \textbf{Low-variance flag.} A student whose own work varies less than
  the cohort's is penalised by z-scoring, and the result says so.
\item
  \textbf{Cohort saturation flag.} When the pooled spread collapses,
  every z-score is a real difference divided by an invented floor, and
  the layer has stopped discriminating. The result says that too.
\end{enumerate}

Mitigations 2, 3, 4 and 5 all have the same character: the system
reports that it cannot tell, instead of producing a number that looks
like an answer. That is the behaviour an evidence system owes a
reviewer.

\hypertarget{layer-3-specified-justified-and-not-built}{%
\section{Layer 3 --- specified, justified, and not
built}\label{layer-3-specified-justified-and-not-built}}

\textbf{Layer 3 is not implemented.} This section states the decision
and its reasoning, because an unexplained hole where a third of the
architecture should be is worse than a documented exclusion.

The layer would have detected machine-generated code, following the
CodeBERT fine-tuning approach of Nguyen et
al.~\cite{nguyen2023, gptsniffer2024}. That work reports two findings
that together make the layer unbuildable \emph{as evidence} in this
setting: classifier performance improves when the training data
resembles the target code, and when snippets are classified in pairs.

A classifier whose accuracy depends on in-distribution training data
requires a labelled corpus of student submissions from the course it
will police. No such corpus exists for this project, and one cannot be
manufactured honestly --- generating synthetic ``AI-written''
submissions and treating them as ground truth would produce a figure
that measures the generator, not the classifier.

The layer could still have been \emph{implemented}. What it could not be
is \emph{characterised}. In a system whose entire premise is auditable
evidence, a layer that emits a number nobody can put a confidence
interval on is worse than an absent layer, because a reviewer has no
basis on which to weigh it against the layers that are measured.

The decision is recorded in the data rather than in a comment. Every
detection result stores:

\begin{verbatim}
aiContent: { status: "skipped", reason: "Layer 3 is not implemented yet" }
\end{verbatim}

so that a result read years later says why it has two layers and not
three. Chapter 8 lists the layer as future work with the labelled corpus
named as its precondition.

\hypertarget{composing-the-review-priority-score}{%
\section{Composing the Review Priority
Score}\label{composing-the-review-priority-score}}

\begin{verbatim}
activeWeight = (structuralOk ? w1 : 0) + (behaviouralOk ? effectiveW2 : 0)
RPS          = (w1 * structural + effectiveW2 * behavioural) / activeWeight
\end{verbatim}

\texttt{services/detection.ts} performs this composition. Three
properties of it are worth a paragraph each in the report, and Figure 5
illustrates them.

\textbf{The configured weights are never the applied ones.} The
configuration carries 0.4 / 0.3 / 0.3 across three layers. With Layer 3
absent the divisor is the sum of the weights of the layers that ran,
which is 0.7, so the applied weights are \textbf{0.571 and 0.429}.
Quoting 0.4 / 0.3 / 0.3 as the operative weights would be wrong, and the
evidence page renders the renormalisation explicitly for this reason.

\textbf{An absent layer contributes nothing, not zero.} The distinction
is not pedantry. A zero is a finding --- it asserts that the layer
looked and found nothing. An absent layer has not looked. Treating
absence as zero would drag every score down by a third and make a
two-layer result incomparable with a three-layer one.

\textbf{The weights and threshold are stored on the result.}
\texttt{effectiveW2} is Layer 2's weight after attenuation by baseline
confidence, and both it and the review threshold are written onto each
\texttt{DetectionResult}. A result therefore remains readable after the
configuration changes, which is what makes the evidence view
reproducible from stored data instead of recomputed against whatever the
configuration happens to be when someone opens it.

\hypertarget{the-queue-and-why-detection-is-not-synchronous}{%
\section{The queue, and why detection is not
synchronous}\label{the-queue-and-why-detection-is-not-synchronous}}

Detection is slow: it parses, it compares against candidates, and it
waits on a second service. Doing it inside the upload request would make
an upload take seconds and fail under load at precisely the moment every
student submits.

\texttt{POST\ /submissions} therefore stores the file, records the
submission and enqueues the work on Redis, \textbf{debounced by job
identity} (Figure 2, step 4). Sixty uploads arriving at a deadline
collapse into a bounded number of detection runs rather than sixty,
because each run reads the current cohort anyway.

A re-run supersedes rather than replaces: \texttt{DetectionResult}
carries \texttt{isCurrent} and \texttt{revision}, so the superseded
result survives. A decision may already have been recorded against it,
and destroying the result would orphan the decision.

\hypertarget{the-web-interface}{%
\section{The web interface}\label{the-web-interface}}

Ten pages and five shared components in \texttt{apps/web}, built between
14 and 26 October. The three worth discussing in the report are:

\begin{itemize}
\tightlist
\item
  \textbf{The review queue} --- assignments ranked by RPS descending,
  showing each student's name and roll number. The roll number is not
  decoration: a fixture containing two students sharing a surname made
  it necessary.
\item
  \textbf{The evidence page} --- the system's actual thesis made
  visible. It renders the caveats \emph{before} the number, the
  per-feature breakdown, the matched spans, and the renormalisation
  arithmetic. It also \textbf{guards its own equation}: it recomputes
  the RPS from the stored weights and, where the recomputation and the
  stored score disagree by more than a small tolerance, refuses to
  render the equation and says the two disagree, rather than printing an
  identity that does not hold.
\item
  \textbf{The student's own result view} --- status, any recorded
  decision, and a form to answer it. \textbf{No score appears anywhere
  on it}, because the endpoint it reads returns none. A ranking a
  student can read is a ranking a student can tune against.
\end{itemize}

A methodology note belongs in Chapter 6 rather than here, but it
originates in this work: rendering real data through these screens found
defects the test suite had not, and printing them to paper found more. A
test asserts a value; these were things nobody had thought to assert.

\hypertarget{what-is-deliberately-not-built}{%
\section{What is deliberately not
built}\label{what-is-deliberately-not-built}}

Stated here so that the presentation's architecture slides and the
delivered system can be reconciled rather than silently differing:

\begin{itemize}
\tightlist
\item
  \textbf{Layer 3}, for the reasons in \S{}5.4.
\item
  \textbf{WebSocket push} for queue updates. The queue is polled by the
  reviewer.
\item
  \textbf{Side-by-side AST diff} rendering of two sources. The evidence
  page reports matching line ranges and links to each file; rendering
  both sources against each other needs a diff view that was not built.
\item
  \textbf{Multi-file and archive submissions.} One source file per
  submission.
\end{itemize}
